% Modern CV, Mesfin Genie
% Compile with: pdflatex mesfin-genie-cv.tex   (run twice)
% Single file, no external class, no system fonts. Compiles on any TeX Live
% or MiKTeX install and on Overleaf without changing a setting.

\documentclass[11pt,a4paper]{article}

\usepackage[T1]{fontenc}
\usepackage[utf8]{inputenc}
\usepackage{lmodern}
\usepackage[margin=2.0cm,top=1.8cm,bottom=1.8cm]{geometry}
\usepackage{xcolor}
\usepackage{titlesec}
\usepackage{enumitem}
\usepackage{tabularx}
\usepackage[hidelinks]{hyperref}
\usepackage{fancyhdr}

\definecolor{ink}{HTML}{16191D}
\definecolor{muted}{HTML}{6A7178}
\definecolor{accent}{HTML}{9A1F35}
\definecolor{rule}{HTML}{E3E5E8}

% Sans throughout, which reads as contemporary rather than as a journal article.
\renewcommand{\familydefault}{\sfdefault}

\color{ink}
\setlength{\parindent}{0pt}
\linespread{1.06}

% Section headings: small, spaced, with a hairline under
\titleformat{\section}
  {\normalfont\large\bfseries\color{ink}}{}{0em}{}[\vspace{-0.55em}{\color{rule}\rule{\linewidth}{0.5pt}}]
\titlespacing*{\section}{0pt}{1.35em}{0.55em}

\newcommand{\entry}[3]{%
  \begin{tabularx}{\linewidth}{@{}p{3.2cm}X@{}}
  \textcolor{muted}{\small #1} & \textbf{#2}\newline #3 \\
  \end{tabularx}\vspace{0.45em}}

\newcommand{\simpleentry}[2]{%
  \begin{tabularx}{\linewidth}{@{}p{3.2cm}X@{}}
  \textcolor{muted}{\small #1} & #2 \\
  \end{tabularx}\vspace{0.35em}}

\pagestyle{fancy}\fancyhf{}
\renewcommand{\headrulewidth}{0pt}
\fancyfoot[C]{\textcolor{muted}{\small Mesfin Genie \quad mesfingenie.com \quad \thepage}}

\begin{document}

{\fontsize{26}{30}\selectfont\bfseries Mesfin Genie}\\[0.35em]
{\color{muted}Senior Lecturer in Health Economics, Newcastle Business School, The University of Newcastle}\\[0.5em]
{\small
\href{mailto:mesfin.genie@newcastle.edu.au}{mesfin.genie@newcastle.edu.au} \quad
\href{https://mesfingenie.com}{mesfingenie.com} \quad
\href{https://orcid.org/0000-0002-1744-4666}{ORCID 0000-0002-1744-4666} \quad
\href{https://scholar.google.com/citations?user=v2SXM_kAAAAJ}{Google Scholar}}

\vspace{0.9em}
{\color{accent}\rule{\linewidth}{1.4pt}}
\vspace{0.4em}

\section*{Profile}
Health economist working on preference measurement: what people are willing to
trade in healthcare, how reliably it can be measured, and what happens to the
measurement when respondents have not read everything put in front of them.
Combines discrete choice experiments with eye tracking and applied econometrics,
and builds the results into decision-support tools used by policymakers.
44 publications across the Journal of Health Economics, Health Economics,
Social Science \& Medicine, Value in Health and npj Vaccines.

\section*{Appointments}
\entry{2026--}{Senior Lecturer in Health Economics}{Newcastle Business School, The University of Newcastle, Australia. Continuing.}
\entry{2023--2025}{Lecturer in Health Economics}{Newcastle Business School, The University of Newcastle, Australia.}
\entry{2023--2024}{Health Measurement and Regulatory Science Fellow}{School of Medicine, Duke University, United States.}
\entry{2022--2023}{Postdoctoral Research Fellow}{Harrison School of Pharmacy, Auburn University, United States.}
\entry{2020--2022}{Research Fellow}{Health Economics Research Unit, University of Aberdeen, United Kingdom.}
\entry{2019--2020}{Postdoctoral Research Fellow}{Department of Economics, Ca' Foscari University of Venice, Italy.}

\section*{Education}
\entry{2019}{PhD, Economics}{Ca' Foscari University of Venice. Stated-preference methods for healthcare decision-making.}
\entry{2015}{MSc, Health Economics and Management}{University of Bologna.}
\simpleentry{}{\textbf{MSc, Economics} \newline Addis Ababa University.}
\simpleentry{}{\textbf{BA, Economics} \newline Jimma University.}

\section*{Selected publications}
\begin{enumerate}[leftmargin=1.1em,itemsep=0.35em,label=\arabic*.]
\item Genie MG, Watson V, Schwarzinger M, Luchini S (2026). Fit for purpose? Assessing the robustness of discrete choice experiment designs in the context of evolving health technologies. \textit{Social Science \& Medicine} 404, 119463.
\item Genie MG, Reed SD, Ozdemir S (2026). Guidance or misdirection? Unpacking the role of feedback in health preference assessments. \textit{Health Economics} 35, 910--928.
\item Martinenghi FI, \ldots\ Genie M, Attwell K, Blyth CC (2026). An open repository of COVID-19 vaccine mandate studies with a worked scoping review of quasi-experimental evidence. \textit{npj Vaccines} 11, 01454.
\item Genie MG, Ngorsuraches S (2026). Priority for self or others? Incorporating equity considerations in a preference-based health value assessment. \textit{Value in Health} 29, 467--475.
\item Genie MG, Boyes AW, Paolucci F (2026). Feeling lonely? Preferences for support programmes to reduce loneliness among older adults in Australia. \textit{Health Policy} 167, 105587.
\item Attema AE, Antonini M, Genie M, Torbica A, Paolucci F (2026). Time preferences and COVID-19 vaccination uptake. \textit{European Journal of Health Economics} 27, 47--63.
\item Genie MG, Poudel N, Paolucci F, Ngorsuraches S (2024). Choice consistency in discrete choice experiments: does numeracy skill matter? \textit{Value in Health} 27, 1594--1604.
\item Genie MG, Ryan M, Krucien N (2023). Keeping an eye on cost: what can eye tracking tell us about attention to cost information in discrete choice experiments? \textit{Health Economics} 32, 1101--1119.
\item Genie MG, Krucien N, Ryan M (2021). Weighting or aggregating? Investigating information processing in multi-attribute choices. \textit{Health Economics} 30, 1291--1305.
\item Genie MG, Nicol\`{o} A, Pasini G (2020). The role of heterogeneity of patients' preferences in kidney transplantation. \textit{Journal of Health Economics} 72, 102331.
\end{enumerate}
\vspace{0.2em}
{\small\color{muted}Full list of 44 items at \href{https://mesfingenie.com/publications/}{mesfingenie.com/publications}.}

\section*{Grants}
\entry{2023--2025}{MandEval, Medical Research Future Fund}{Effectiveness and consequences of Australia's COVID-19 vaccine mandates. MRF2019107. Programme A\$4,754,183; University of Newcastle allocation A\$808,440. Investigator.}
\entry{2025}{World Bank Group, IBRD}{Business case for upscaling Field Epidemiology Training Programmes in India. \$416,000.}
\simpleentry{2019}{American Society of Health Economists travel award, with the NIA-funded Center on the Demography and Economics of Health and Aging, Stanford University.}

\section*{Decision-support tools}
\simpleentry{eMANDEVAL}{Comparing vaccine mandate designs on predicted public support, health outcomes and cost.}
\simpleentry{STEPS}{Planning Field Epidemiology Training Programme scale-up in India, with the World Bank.}
\simpleentry{SOIL CRC BCA}{Benefit-cost appraisal for soil and farming interventions. Project lead: Dr Frank Wogbe Agbola.}

\section*{Teaching}
\simpleentry{ECON3111}{Behavioural Economics. Lecturer and Course Coordinator.}
\simpleentry{ECON2112}{Health Economics and Finance. Lecturer and Course Coordinator.}
\simpleentry{ECON3011}{Introduction to Health Economics and Finance. Primary Course Coordinator.}

\section*{Supervision}
Four current PhD candidates, on telehealth use in Australia, climate change and
agriculture in the Solomon Islands, crop drought resilience, and menopause and
female labour market participation. One PhD completed 2025 on health systems
reform, equity and efficiency. Two MPhil students completed on the impacts of
COVID-19 and influenza vaccine mandates.

\section*{Service}
Guest Editor, \textit{Health Policy and Technology}. Reviewer for the
\textit{Journal of Health Economics}, \textit{Health Economics},
\textit{Value in Health}, \textit{Social Science \& Medicine},
\textit{Journal of Choice Modelling}, \textit{Theory and Decision} and
\textit{BMJ Open}.

\end{document}
